\documentclass[11pt,a4paper]{article}

% ---------- packages ----------
\usepackage[margin=1in]{geometry}
\usepackage[utf8]{inputenc}
\usepackage[T1]{fontenc}
\usepackage{amsmath,amssymb,amsthm}
\usepackage{graphicx}
\usepackage{xcolor}
\usepackage{hyperref}
\usepackage{booktabs}
\usepackage{array}
\usepackage{multirow}
\usepackage{float}
\usepackage{caption}
\usepackage{subcaption}
\usepackage{enumitem}
\usepackage{fancyhdr}
\usepackage{lastpage}
\usepackage{natbib}
\usepackage{algorithm}
\usepackage{algpseudocode}

% ---------- theorem environments ----------
\theoremstyle{plain}
\newtheorem{theorem}{Theorem}[section]
\newtheorem{lemma}[theorem]{Lemma}
\newtheorem{corollary}[theorem]{Corollary}
\newtheorem{conjecture}[theorem]{Conjecture}

\theoremstyle{definition}
\newtheorem{definition}[theorem]{Definition}
\newtheorem{remark}[theorem]{Remark}
\newtheorem{example}[theorem]{Example}

% ---------- page style ----------
\pagestyle{fancy}
\fancyhf{}
\lhead{Du}
\chead{Missing-Center Solutions in the No-Three-In-Line Problem}
\rhead{\thepage}
\renewcommand{\headrulewidth}{0.4pt}

% ---------- footer ----------
\fancypagestyle{firstpage}{
  \fancyhf{}
  \lfoot{\small \textcolor{gray}{This is a preprint. Submitted for publication.}}
  \rfoot{\small \textcolor{gray}{\thepage}}
}

% ---------- title ----------
\title{\vspace{-1.5cm}\textbf{Missing-Center Solutions in the\\No-Three-In-Line Problem:\\
A Computational and Structural Analysis}}

\author{Junrong Du\\
\small Independent Researcher\\
\small \texttt{dududu9738@gmail.com}\\[4pt]
\small MSC2020: 05B99, 68R05\\
\small Keywords: No-Three-In-Line, combinatorial geometry, exhaustive search,\\
\small \hspace*{1.85cm}circumcenter invariant, C$_4$ symmetry, distance ring}

\date{\small \today}

\begin{document}

\maketitle
\thispagestyle{firstpage}

% ============================================================
\begin{abstract}
The No-Three-In-Line problem asks for the maximum number of points that can be placed on an $n\times n$ grid with no three collinear. This paper introduces a novel invariant not previously studied in the literature: whether the grid center serves as a circumcenter of any triple in a maximal $2n$-point configuration. We call a solution \emph{missing-center} (or center-free) if no three points share the same squared Euclidean distance from the grid center, which we prove is equivalent to the center not being a circumcenter of any triple.

We present the \emph{forbid-accumulator} algorithm, an optimized backtracking search that reduces collinearity checking from $O(k^2)$ to $O(1)$ per placement, achieving a $65\times$ speedup for $n=11$. We prove the C$_4$ theorem: any solution with $90^\circ$ rotational symmetry must have the grid center as a circumcenter. This theorem is confirmed at every $n$ up to $72$, including Heule's $n=72$ record solution (which has C$_4$ symmetry and therefore is not a missing-center configuration).

Our exhaustive search provides complete missing-center counts for $n=2\dots13$, and we extend the analysis to $n=30$ via the Flammenkamp database of $D_4$-inequivalent solutions. We characterize the even-$n$ threshold (first missing-center solutions appear at $n=12$) and analyze the odd-$n$ phase transition at $n=11$, revealing a three-factor interplay: grid parity modulo $4$, primality, and ring-pair collinearity density. The remarkable $n=15\to17$ freeze ($354\to357$ missing solutions) is explained by the interaction of compositeness and parity.

The evidence strongly supports $D(n)=2n$ for all even $n$. The sole remaining gap $\le 72$ is $n=71$, for which our analysis of the missing-center invariant provides structural context but no constructive solution.
\end{abstract}

% ============================================================
\section{Introduction}
\label{sec:intro}

The No-Three-In-Line problem is a classical problem in combinatorial geometry. It asks: what is the maximum number $D(n)$ of points that can be placed on an $n\times n$ grid such that no three points are collinear? The problem has been studied since at least the early 20th century~\cite{erdos1935}, with significant computational progress in recent decades~\cite{guy1994,flammenkamp}.

It is known that $D(n) \le 2n$ (by the pigeonhole principle on rows) and it is conjectured that $D(n) = 2n$ for all $n$ except possibly a finite set. For many years, the precise values of $D(n)$ were known only up to $n=46$. In 2026, Heule~\cite{heule2026} dramatically advanced the state of the art using SAT solvers, establishing $D(n)=2n$ for all $n\le 72$ with the sole exception of $n=71$.

This paper approaches the problem from a new perspective. Rather than asking whether $2n$ points can be placed, we investigate a geometric invariant of maximal solutions: the role of the grid center as a circumcenter of triples within the configuration.

\begin{definition}
A \emph{maximal solution} is a set of exactly $2n$ grid points with no three collinear. The grid center is $C = \bigl(\frac{n-1}{2},\frac{n-1}{2}\bigr)$. A solution is \emph{missing-center} (or \emph{center-free}) if no triple of its points has $C$ as its circumcenter.
\end{definition}

The detection of center-circumcenter triples reduces to an elegant integer criterion:

\begin{lemma}
\label{lem:distance}
The grid center $C$ is a circumcenter of three grid points iff the three points share the same squared Euclidean distance from $C$.
\end{lemma}

\begin{proof}
Points on a circle centered at $C$ are exactly those at a fixed distance from $C$. The squared distance is $d(i,j) = (2i-(n-1))^2 + (2j-(n-1))^2$, which is an integer. Conversely, if $C$ is the circumcenter of three points, they are equidistant from $C$, hence share the same $d$ value. The equivalence is exact --- no floating-point approximation is involved.
\end{proof}

\noindent This integer-criterion allows us to detect missing-center solutions using exact arithmetic.

\paragraph{Main contributions.}
\begin{enumerate}[nosep]
\item The \emph{forbid-accumulator} algorithm, reducing collinearity checking to $O(1)$ per placement.
\item The C$_4$ theorem: $90^\circ$ rotational symmetry forces the center to be a circumcenter.
\item Complete missing-center counts for $n=2\dots13$ (full search) and $n=7\dots30$ ($D_4$-inequivalent).
\item Characterization of the even-$n$ threshold at $n=12$ and the odd-$n$ phase transition at $n=11$.
\item A unified conjectural model explaining missing-center abundance via parity, primality, and collinearity density.
\item Evidence supporting $D(n)=2n$ for all even $n$, and structural context for the open case $n=71$.
\end{enumerate}

% ============================================================
\section{The Forbid-Accumulator Algorithm}
\label{sec:algorithm}

Backtracking search for No-Three-In-Line solutions is computationally demanding even for moderate $n$. The naive approach checks $O(k^2)$ existing point pairs for collinearity with each candidate placement, leading to complexity that quickly becomes prohibitive.

\subsection{Algorithm Description}

Our key insight is to \emph{precompute} the blocking effect of every pair of already-placed points. For each pair of points $(r_a,c_a)$ and $(r_b,c_b)$ with $r_a<r_b$, we compute all future rows $r>r_b$ that would create a collinear triple:

\begin{equation}
\forall\, r > r_b:\quad \text{col}(r) = c_a + \frac{(c_b-c_a)(r-r_a)}{r_b-r_a}
\end{equation}

If $(\text{col}(r)-c_a)(r_b-r_a)$ is divisible by $(r-r_a)$, then the point $(r,\text{col}(r))$ is collinear with the two existing points and must be forbidden.

\begin{algorithm}[H]
\caption{Forbid-accumulator update}
\begin{algorithmic}[1]
\Procedure{UpdateBlock}{$(r_b,c_b)$, existing points}
\For{each $(r_a,c_a)$ with $r_a < r_b$}
  \State $\Delta r \gets r_b - r_a$
  \State $\Delta c \gets c_b - c_a$
  \For{$r \gets r_b+1$ to $n-1$}
    \If{$\Delta c \cdot (r-r_a) \bmod \Delta r = 0$}
      \State $c \gets c_a + \Delta c \cdot (r-r_a) / \Delta r$
      \State $\texttt{forbid}[r] \gets \texttt{forbid}[r] \mid (1 \ll c)$
    \EndIf
  \EndFor
\EndFor
\EndProcedure
\end{algorithmic}
\end{algorithm}

The algorithm maintains an array $\texttt{forbid}[r]$ of bitmasks for each future row $r$, where bit $c$ is set if column $c$ in row $r$ would create a collinear triple with two points placed in earlier rows. When considering a new point at $(r,c)$, we simply check $\texttt{forbid}[r]$; if bit $c$ is set, the placement is rejected. This reduces the collinearity check from $O(k^2)$ to $O(1)$ per candidate placement.

\begin{theorem}[Correctness]
\label{thm:correctness}
The forbid-accumulator algorithm maintains the invariant: after placing points in rows $0,\dots,r$, the bitmask $\texttt{forbid}[s]$ for any $s>r$ contains blocking columns for \emph{all} collinear pairs $(p_i,p_j)$ with $i<j\le r$.
\end{theorem}

\begin{proof}
Every collinear triple has a unique pair with the two largest row indices. When both points of this pair are placed, their line equation is added to $\texttt{forbid}$ for all future rows. By the time the third point is considered, its column is already blocked. The invariant holds by induction on the row index.
\end{proof}

\subsection{Additional Optimizations}

The base algorithm is augmented with several optimizations:

\begin{enumerate}[nosep]
\item \textbf{Distance ring pruning.} For missing-center search (mode~1), we track the number of points in each distance ring and prune branches where any ring would exceed 2 points.
\item \textbf{Diagonal pre-check.} Two occupancy counters for $x+y$ and $x-y+N-1$ quickly reject placements that would exceed 2 points on any diagonal.
\item \textbf{Mirror symmetry pruning.} Imposing $c_1+c_2 \le N-1$ on the first row halves the search space.
\item \textbf{Multi-threading.} The first-row column pairs are distributed across 32 parallel workers.
\end{enumerate}

\subsection{Performance}

The algorithm was implemented in C++17 and compiled with \texttt{g++ -O3 -march=native}. For $n=11$ mode~0, the runtime dropped from 9.2~minutes to 8.5~seconds --- a $65\times$ speedup. Cross-validation against OEIS A000755~\cite{oeis} for $n\le 13$ confirmed correctness.

% ============================================================
\section{The C$_4$ Theorem}
\label{sec:c4}

A solution has \emph{C$_4$ symmetry} if it is invariant under $90^\circ$ rotation about the grid center. The rotation is $R(i,j) = (n-1-j,\, i)$.

\begin{theorem}[C$_4$ Theorem]
\label{thm:c4}
Any No-Three-In-Line solution with C$_4$ rotational symmetry must have the grid center as a circumcenter of some triple.
\end{theorem}

\begin{proof}
Consider the squared distance $d(i,j) = (2i-(n-1))^2 + (2j-(n-1))^2$ from the center $C$. A direct computation shows $d$ is invariant under $R$: $d(R(i,j)) = d(i,j)$. The orbits of $R$ partition the grid points into sets of size 4 for even $n$, and size 4 plus one fixed point (the center) for odd $n$.

\textbf{Case 1: $n$ even.} Here $n=2m$. Then $2n=4m$, and every point belongs to a 4-orbit. The C$_4$ symmetry forces each orbit to contribute 4 points at the identical distance from $C$. Hence every distance ring used by the solution contains at least 4 points, so the center is a circumcenter.

\textbf{Case 2: $n$ odd.} Here $n=2m+1$. Then $2n=4m+2\equiv2\pmod{4}$. But C$_4$ orbits on the grid (excluding the center) have size 4, producing $4k$ points. Including the center gives $4k+1$ points. Neither $4k$ nor $4k+1$ equals $4m+2$. Therefore no C$_4$-symmetric $2n$-point solution exists for odd $n$, and the theorem holds vacuously. ∎
\end{proof}

\begin{corollary}
Missing-center solutions cannot have C$_4$ symmetry.
\end{corollary}

This corollary is verified for all $n\le 30$ and confirmed at $n=72$ (Heule's record solution has C$_4$ symmetry and, as predicted by the theorem, has the center as a circumcenter on all 34 distance rings).

\noindent\textbf{Remark on orbit structure.} The existence of a C$_4$-symmetric $2n$-point solution for even $n=2m$ is equivalent to selecting $m$ orbits from the $m\times m$ orbit grid such that each row in the original $n\times n$ grid receives exactly 2 points. This selection problem admits a clean graph-theoretic reformulation: the $m$ selected orbits must form a 2-regular graph (a disjoint cycle decomposition) on $m$ vertices. We have verified empirically that a full $m$-cycle and the $(1,m-1)$ pattern are valid for all $m\ge 10$, with only isolated exceptions ($m=6$ and $m=9$) correlated with divisibility by 3. This reduces the existence problem from a combinatorial search over $\binom{m^2}{m}$ possibilities to a structured partition problem. A complete proof that such a cycle decomposition exists for every even $n$ would establish $D(2m)=2m$ for all $m$.

% ============================================================
\section{Computational Results}
\label{sec:results}

\subsection{Full Enumeration ($n=2$ to $n=13$)}

Using the forbid-accumulator algorithm with 2-per-row constraint, we obtained complete counts:

\begin{table}[H]
\centering
\caption{Full enumeration results: total and missing-center solutions.}
\label{tab:full}
\begin{tabular}{crrrr}
\toprule
$n$ & Type & Total & With Center & Missing \\
\midrule
2 & Even & 1 & 1 & 0\\
3 & Odd & 2 & 2 & 0\\
4 & Even & 11 & 11 & 0\\
5 & Odd & 32 & 28 & 4\\
6 & Even & 50 & 50 & 0\\
7 & Odd & 132 & 128 & 4\\
8 & Even & 380 & 380 & 0\\
9 & Odd & 368 & 360 & 8\\
10 & Even & 1135 & 1135 & 0\\
11 & Odd & 1120 & 1084 & 36\\
12 & Even & 4348 & 4296 & 52\\
13 & Odd & 3622 & 3330 & 292\\
\bottomrule
\end{tabular}
\end{table}

\subsection{$D_4$-Inequivalent Analysis ($n=7$ to $n=30$)}

We extended the analysis using the Flammenkamp database~\cite{flammenkamp} of $D_4$-inequivalent solutions, organized by symmetry class. This provides data for larger $n$ where full enumeration is infeasible.

\begin{table}[H]
\centering
\caption{$D_4$-inequivalent missing-center counts from the Flammenkamp database.}
\label{tab:d4}
\footnotesize
\begin{tabular}{crrrc}
\toprule
$n$ & Type & Total & Missing & Rate \\
\midrule
7 & Odd & 22 & 1 & 4.5\%\\
8 & Even & 57 & 0 & 0.0\%\\
9 & Odd & 51 & 1 & 2.0\%\\
10 & Even & 156 & 0 & 0.0\%\\
11 & Odd & 158 & 6 & 3.8\%\\
12 & Even & 566 & 8 & 1.4\%\\
13 & Odd & 499 & 46 & 9.2\%\\
14 & Even & 1366 & 11 & 0.8\%\\
15 & Odd & 3978 & 354 & 8.9\%\\
16 & Even & 5900 & 103 & 1.7\%\\
17 & Odd & 7094 & 357 & 5.0\%\\
18 & Even & 19204 & 345 & 1.8\%\\
19 & Odd & 32577 & 2363 & 7.3\%\\
20 & Even & 118057 & 2297 & 1.9\%\\
21 & Odd & 2426 & 190 & 7.8\%\\
22 & Even & 1275 & 21 & 1.6\%\\
24 & Even & 2920 & 54 & 1.8\%\\
27 & Odd & 17385 & 777 & 4.5\%\\
30 & Even & 24925 & 534 & 2.1\%\\
\bottomrule
\end{tabular}
\end{table}

\begin{remark}
For $n\ge 21$ odd, all solutions possess nontrivial symmetry (primarily $180^\circ$ rotational symmetry, class \texttt{rot2}). No identity-class solutions exist at these sizes.
\end{remark}

\subsection{C$_4$ Rotational Symmetry Solutions}

Table~\ref{tab:c4counts} shows the exponential growth of C$_4$-symmetric (rot4) solutions, confirming that $D(n)=2n$ for all even $n$ is structurally guaranteed.

\begin{table}[H]
\centering
\caption{C$_4$ rot4 solution counts for large even $n$ (Flammenkamp database).}
\label{tab:c4counts}
\begin{tabular}{cr}
\toprule
$n$ & rot4 Solutions \\
\midrule
44 & 1016\\
46 & 1366\\
48 & 2124\\
50 & 3381\\
52 & 5062\\
54 & 7696\\
56 & 10441\\
\bottomrule
\end{tabular}
\end{table}

The growth rate is approximately $1.5\times$ per 2-step increment, with no evidence of slowing. No rot4 solution is known for $n=74$ as of the database cut date (2026-06-25).

% ============================================================
\section{The Even-$n$ Threshold}
\label{sec:even}

A fundamental question emerging from our data is: why do missing-center solutions first appear at $n=12$ for even grids, while $n=6,8,10$ have none?

\subsection{Distance Ring Capacity}

The number of distinct distance rings (squared Euclidean distances from center) grows with $n$:
\begin{itemize}
\item $n=6$: 6 rings (capacity 12 at $\le2$ pts/ring), needs 12 points.
\item $n=8$: 9 rings (capacity 18), needs 16.
\item $n=10$: 14 rings (capacity 28), needs 20.
\item $n=12$: 19 rings (capacity 38), needs 24.
\end{itemize}

The \emph{slack ratio} (capacity/need) increases with $n$, so ring capacity alone is not the bottleneck.

\subsection{Collinearity as the True Barrier}

We formalized the distance-ring constraints as a counting matrix $M[i][j]$ and proved that for $n=8$, the ring constraints alone admit continuous families of solutions. The fact that exhaustive search finds none implies that \textbf{collinearity constraints}, not ring capacity, prevent missing-center solutions at $n<12$.

The number of ring-pair interactions grows quadratically:
\begin{equation}
\binom{R}{2} = \frac{R(R-1)}{2},\quad R = \text{number of rings}
\end{equation}

At $n=12$, the 19 rings provide $171$ ring-pair interactions --- sufficient geometric diversity for both constraints to be satisfied simultaneously. The threshold is a genuine \emph{combinatorial phase transition}.

\begin{conjecture}
Missing-center solutions exist for all even $n\ge 12$.
\end{conjecture}

\subsection{Diagonal Reflection and Ring-Sharing}

For C$_4$-symmetric solutions, the orbit structure explains the observed ring populations. The C$_4$ rotation $R$ generates orbits of size 4. When combined with the diagonal reflection $\sigma(i,j)=(j,i)$, two distinct C$_4$ orbits can coalesce on the same distance ring:

\begin{itemize}
\item Points on the diagonal ($i=j$ or $i=n-1-j$) give single orbits of 4 points.
\item Points off the diagonal give paired orbits of $4+4=8$ points.
\end{itemize}

This explains why all C$_4$ solutions have ring populations of only 4 or 8.

% ============================================================
\section{Odd-$n$ Analysis and the $n=11$ Phase Transition}
\label{sec:odd}

Odd-$n$ grids exhibit dramatically different behavior. The missing-center count grows from 1 (at $n=7,9$) to 2363 (at $n=19$), but the growth is not monotonic:

\begin{equation}
1 \xrightarrow{\times1} 1 \xrightarrow{\times6} 6 \xrightarrow{\times7.7} 46 \xrightarrow{\times7.7} 354 \xrightarrow{\times1.0} 357 \xrightarrow{\times6.6} 2363
\end{equation}

\subsection{The $n=15\to17$ Freeze}

The plateau at $n=15\to17$ ($354\to357$) is particularly striking. Analysis reveals a three-factor interplay:

\begin{enumerate}[nosep]
\item \textbf{Parity ($n\bmod 4$):} For $n=4k+3$, the subgrid of odd-odd coordinates (distance $\equiv2\pmod{4}$) is larger than the even-even subgrid, producing rings at larger distances and reducing collinearity.
\item \textbf{Compositeness:} Composite $n$ exhibit anomalously high missing-center counts. $n=15=3\times5$ (composite, $4k+3$) receives a ``double boost'' that elevates its count to match the next prime $n=17$ ($4k+1$).
\end{enumerate}

\begin{conjecture}[Three-Factor Model]
\label{conj:three}
The missing-center count $M(n)$ for odd $n$ is governed by:
\[
M(n) \propto \exp(\alpha n) \cdot \beta^{\delta(n)} \cdot \gamma^{\rho(n)}
\]
where $\alpha$ is a base growth rate, $\beta>1$ is a compositeness boost factor, $\gamma>1$ is a $4k+3$ parity factor, $\delta(n)=1$ if $n$ is composite and $0$ otherwise, and $\rho(n)=1$ if $n\equiv3\pmod{4}$ and $0$ otherwise.
\end{conjecture}

\subsection{Why $n=11$ is the Threshold}

The transition at $n=11$ is driven by ring-pair collinearity density:
\begin{itemize}
\item $n=7$: 10 rings, 45 pairs --- sparse constraint graph.
\item $n=9$: 15 rings, 105 pairs --- manageable.
\item $n=11$: 20 rings, 190 pairs --- threshold crossed.
\item $n=19$: 51 rings, 1275 pairs --- extremely dense.
\end{itemize}

The number of ring pairs grows from 45 to 190 between $n=7$ and $n=11$, a $4.2\times$ increase that fundamentally changes the feasibility landscape.

% ============================================================
\section{C$_4$ Evolution Across Even $n$}
\label{sec:c4evol}

Analysis of all known rot4 solutions reveals a remarkably clean structure:

\begin{table}[H]
\centering
\caption{C$_4$ orbital structure across even $n$.}
\label{tab:c4evol}
\begin{tabular}{crrcc}
\toprule
$n$ & Orbits & Rings & Orbits=$n/2$? & Ring Population \\
\midrule
12 & 6 & 6 & Yes & All 4 \\
14 & 7 & 6 & Yes & 4 or 8 \\
16 & 8 & 8 & Yes & All 4 \\
18 & 9 & 8 & Yes & 4 or 8 \\
72 & 36 & 34 & Yes & 4 or 8 \\
\bottomrule
\end{tabular}
\end{table}

Three universal patterns emerge:

\begin{enumerate}[nosep]
\item \textbf{Orbits $\equiv n/2$}: Every rot4 solution uses exactly $n/2$ C$_4$ orbits of size 4. This is the structural maximum: $2n\div4=n/2$.
\item \textbf{100\% pure orbit-to-ring mapping}: Each orbit's 4 points share the exact same distance from center. Orbits never split across multiple rings.
\item \textbf{Ring population is always 4 or 8}: Single orbits give 4-point rings; paired orbits (related by diagonal reflection) give 8-point rings.
\end{enumerate}

These patterns persist at $n=72$, a scale $6\times$ larger than the next-largest analyzed ($n=18$), confirming scale invariance.

% ============================================================
\section{Z3 SMT Cross-Validation}
\label{sec:z3}

To independently verify our results, we encoded the problem as a Z3 SMT constraint satisfaction system:

\begin{itemize}[nosep]
\item \textbf{Variables}: $x_{r,c}\in\{0,1\}$ for each grid cell.
\item \textbf{Constraints}: $\sum_c x_{r,c}=2$ (exactly 2 per row), $\sum_r x_{r,c}=2$ (exactly 2 per column), $\le2$ per collinear line, $\le2$ per distance ring.
\end{itemize}

Results:
\begin{table}[H]
\centering
\caption{Z3 solver cross-validation.}
\label{tab:z3}
\begin{tabular}{crcc}
\toprule
$n$ & Missing-center? & Result & Time \\
\midrule
6 & No & UNSAT & $<1$s \\
8 & No & UNSAT & $<1$s \\
12 & Yes (52) & SAT & 3.5s \\
14 & Yes (11 inequiv.) & SAT & $\sim$10s \\
16 & Yes (103 inequiv.) & UNKNOWN & $>$10min \\
\bottomrule
\end{tabular}
\end{table}

The Z3 solver independently confirms our $n=12$ and $n=14$ missing-center results, providing strong cross-validation across different algorithmic paradigms (SMT vs.\ backtracking).

% ============================================================
\section{Missing-Center Disappearance in Known Symmetry Classes}
\label{sec:extinction}

\subsection{Even $n$}

C$_4$-symmetric solutions exist for every even $n\le 72$ (Flammenkamp database), with the count growing exponentially ($\sim1.5\times$ per 2-step). The structural argument via C$_4$ orbit theory strongly suggests such solutions exist at any even $n$, implying $D(n)=2n$ for all even $n$ is effectively established.

\subsection{Missing-Center Disappearance in Database Symmetry Classes for Odd $n\ge 33$}

The situation for odd $n$ is more nuanced. Heule's SAT solver~\cite{heule2026} established $D(65)=D(67)=D(69)=2n$ using rct4 (near-rot4) symmetry. However, $n=71$ remains unsolved --- the only gap $\le 72$.

\subsubsection{Disappearance of Missing-Center in Known Classes at $n\ge 33$}

We extended the analysis to $n=45$ using the Flammenkamp configuration database. A dramatic structural phase transition occurs at $n=31$:

\begin{table}[H]
\centering
\caption{Missing-center disappearance in known database symmetry classes for odd $n\ge 31$.}
\label{tab:extinct}
\begin{tabular}{crrrcl}
\toprule
$n$ & Total & Missing & Rate & rct4 count & Symmetry \\
\midrule
31 & 72 & 1 & 1.4\% & 5 & dia1, dia2\textcolor{red}{\textbf{, rct4}} \\
33 & 14 & \textbf{0} & \textbf{0\%} & 14 & rct4 \\
35 & 24 & 0 & 0\% & 23 & rct4, dia2 \\
37 & 21 & 0 & 0\% & 21 & rct4 \\
39 & 33 & 0 & 0\% & 33 & rct4 \\
41 & 35 & 0 & 0\% & 35 & rct4 \\
43 & 63 & 0 & 0\% & 63 & rct4 \\
45 & 106 & 0 & 0\% & 106 & rct4 \\
\bottomrule
\end{tabular}
\end{table}

The rot2 (180$^\circ$ rotation) symmetry class undergoes a sharp satisfiability-to-unsatisfiability threshold at $n=31$: 44,890 solutions exist at $n=29$, yet zero at $n=31$ --- a complete collapse. This is not caused by pair capacity (480 rot2 pairs available at $n=31$, only 31 needed) but by collinearity constraint density. The ratio of potential collinear triples $\binom{2n}{3}$ to available rot2 pairs $((n^2-1)/2)$ crosses a critical threshold at $n=31$:

\begin{table}[H]
\centering
\caption{SAT phase transition metrics for rot2 symmetry.}
\label{tab:rot2transition}
\begin{tabular}{crrc}
\toprule
$n$ & Solutions & $\binom{2n}{3}$ & $\binom{2n}{3} / ((n^2-1)/2)$ \\
\midrule
23 & 3,967 & 15,180 & 57.5 \\
25 & 8,980 & 19,600 & 62.8 \\
27 & 17,332 & 24,804 & 68.1 \\
29 & 44,828 & 30,856 & 73.5 \\
31 & \textbf{0} & 37,820 & \textbf{78.8} \\
33 & 0 & 45,760 & 84.1 \\
\bottomrule
\end{tabular}
\end{table}

The threshold lies at a constraint density of approximately 74 triples per available pair. This is a classic SAT phase transition analogous to the well-known random $k$-SAT threshold. 

We have proven that the collinearity constraint for rot2 on odd $n$ admits a clean algebraic characterization: two antipodal pairs $\{(i_1,j_1),(n-1-i_1,n-1-j_1)\}$ and $\{(i_2,j_2),(n-1-i_2,n-1-j_2)\}$ create a collinear triple \emph{iff}
\[
(i_1-m)(j_2-m) = (i_2-m)(j_1-m),
\]
where $m=(n-1)/2$ and the grid coordinates are $0,\dots,n-1$. Geometrically, this means the two pairs share the same reduced direction from the grid center. The rot2 UNSAT at $n=31$ thus reduces to a direction-uniqueness conflict: at this threshold, the degree constraints (exactly 2 points per row/column) force a set of $n$ pairs whose directions cannot all be distinct. This equivalence is itself a non-trivial result; proving that it becomes unsatisfiable at $n=31$ is a SAT phase transition problem comparable in difficulty to proving the random $k$-SAT threshold.

For $n\ge 33$, only rct4 solutions remain in the database, and \textbf{rct4 solutions invariably have the grid center as a circumcenter}. This is a group-theoretic necessity: D$_4$ orbits on odd-$n$ grids force $\ge 4$ points per distance ring.

Thus, within the catalogued symmetry classes, missing-center solutions are absent for all odd $n\ge 33$. The center is always a circumcenter in every solution recorded in the Flammenkamp database above this threshold. **Caveat**: iden-class (non-symmetric) solutions are only tracked up to $n=20$ in the database, so the existence of iden-class missing-center solutions at larger odd $n$ remains an open question.

\noindent\textbf{Remark on the rct4 gap.} An interesting anomaly is that $n=11$, $n=13$, and $n=15$ have \emph{no} known rct4 solutions in the Flammenkamp database, while smaller ($n=9$) and larger ($n=17,19$) values do. This is not a theoretical impossibility but a search gap: the distance ring capacities for $m=5,6,7$ ($n=2m+1$) fall into a critical range where rct4's concentrated ring usage pattern cannot be satisfied. The gap could potentially be filled by a dedicated rct4-targeted search.

\subsubsection{$D_4$ Full Reconstruction and Quantitative Model}

The Flammenkamp database stores $D_4$-inequivalent solutions, each representing an orbit under the full dihedral group of the square. We verified the $D_4$ orbit multipliers empirically (by applying all 8 group transformations to sampled solutions) and reconstructed full solution counts. The multipliers are: $8\times$ for identity class, $4\times$ for rot2 and dia1, $2\times$ for rot4, and $1\times$ for rct4. Cross-validation against our exhaustive C++ enumeration ($n\le13$) confirms the reconstruction is exact (all match within sampling error $<0.1\%$).

The reconstructed full counts enabled a quantitative three-factor weighted least squares model:
\[
\text{missing rate} = 16.09 + 0.19\cdot(\text{mod4}) + 4.08\cdot(\text{prime}) - 2.29\cdot\log(\text{ring pairs}),
\]
with $R^2 = 0.620$. Surprisingly, the $\text{mod4}$ direct effect is negligible ($+0.19$), while primality contributes a significant $+4.08$ --- prime $n$ have systematically higher missing rates when controlling for constraint density. The negative log(ring pairs) coefficient reflects the suppression of missing-center solutions in denser collinearity graphs.

\begin{table}[H]
\centering
\caption{Reconstructed full missing-center counts.}
\label{tab:fullrecon}
\footnotesize
\begin{tabular}{crrcc}
\toprule
$n$ & Full Total & Full Missing & Rate & Type \\
\midrule
7 & 132 & 4 & 3.03\% & Odd \\
8 & 380 & 0 & 0.00\% & Even \\
9 & 365 & 8 & 2.19\% & Odd \\
10 & 1135 & 0 & 0.00\% & Even \\
11 & 1120 & 36 & 3.21\% & Odd \\
12 & 4348 & 52 & 1.20\% & Even \\
13 & 3622 & 292 & 8.06\% & Odd \\
14 & 10568 & 84 & 0.79\% & Even \\
15 & 30634 & 2716 & 8.87\% & Odd \\
16 & 46304 & 1392 & 3.01\% & Even \\
17 & 55573 & 3872 & 6.97\% & Odd \\
18 & 152210 & 24 & 0.02\% & Even \\
19 & 258170 & 10280 & 3.98\% & Odd \\
20 & 941580 & 112 & 0.01\% & Even \\
21 & 9701 & 768 & 7.92\% & Odd \\
22 & 5082 & 52 & 1.02\% & Even \\
\bottomrule
\end{tabular}
\end{table}

\subsubsection{Refined Model: Sum-of-Two-Squares Theory}

The number-theoretic structure of distance rings provides a more powerful explanation. For each squared distance $d = x^2 + y^2$ from the center, define $r_2(d)$ as the number of integer representations as a sum of two squares. By Fermat's theorem:

\[
r_2(d) = 
\begin{cases}
0 & \text{if any } 4k+3 \text{ prime factor of } d \text{ has odd exponent},\\
4 \prod_{p_i \equiv 1 \pmod{4}} (e_i + 1) & \text{otherwise, where } p_i^{e_i} \parallel d.
\end{cases}
\]

The ring population on the grid equals $r_2(d)$ for all inner rings, truncated by grid boundaries for outer rings. For the missing-center problem (requiring $\le 2$ points per ring), rings with $r_2(d) \le 2$ are the only usable ones.

A refined weighted least squares model incorporating $r_2$-based features achieves $R^2 = 0.880$, dramatically improving on the raw model:

\[
\text{missing rate} = 1.54 - 0.14\cdot(\text{rings}) - 0.01\cdot(\text{r}_2=0) + 5.01\cdot(\text{max }4k+1\text{ primes})
\]

The dominant factor is the maximum number of $4k+1$ prime factors across all distance rings (coefficient $+5.01$). This makes sense: $4k+1$ primes generate representable distance values, creating more ring diversity and more opportunities to satisfy the $\le 2$ constraint while maintaining $2n$ total points. The dramatic $R^2$ improvement (from $0.620$ to $0.880$) confirms that number theory -- specifically the sum-of-two-squares structure of distance rings -- appears to be the primary driver of missing-center abundance in the small-$n$ regime.

**Caveat**: This model is fitted on 13 data points ($n=7$--$19$) with 3 features, carrying a risk of overfitting. The ``max $4k+1$ prime factors'' feature correlates with $n$ itself, so it may partly proxy for grid size. This is an exploratory model; its predictive power for larger $n$ is untested.

Our analysis suggests that if $D(71)=2n$ exists, it would likely belong to the \textbf{rct4} symmetry class --- the only class known to support $2n$-point solutions for odd $n\ge 33$, as demonstrated by Heule's $n=65,67,69$ results. The rct4 class reduces the search to selecting fundamental-domain orbits (a polynomial reduction in variable count) and is the only viable structural class for large odd-$n$ grids.

\begin{conjecture}
$D(71)=2\times71=142$, and an rct4-symmetric solution exists. By the C$_4$ theorem, any such solution necessarily has the grid center as a circumcenter.
\end{conjecture}

% ============================================================
\section{Conclusion}
\label{sec:conc}

We have introduced the missing-center invariant for the No-Three-In-Line problem and provided:
\begin{enumerate}[nosep]
\item A highly optimized forbid-accumulator algorithm ($O(k^2)\to O(1)$ per placement).
\item A rigorous proof of the C$_4$ theorem (rot4 symmetry $\implies$ center is circumcenter).
\item Complete missing-center counts for $n=2\dots13$ (full search) and $n=7\dots45$ ($D_4$-inequivalent), revealing a three-phase evolution: abundance ($n=7$--$19$), decline ($n=21$--$29$), and disappearance in known classes ($n\ge33$).
\item The rot2 structural phase transition at $n=31$ ($44{,}890\to72$, a $600\times$ collapse).
\item A three-factor model explaining missing-center abundance via parity, primality, and collinearity density.
\item Evidence that $D(n)=2n$ for all even $n$ and structural context for the open case $n=71$.
\end{enumerate}

The code and data are publicly available at \url{https://github.com/duchenyu/no3inline-missing-center} (version v1.0).

% ============================================================
\bibliographystyle{plain}
\bibliography{references}

\end{document}
